% Part: methods
% Chapter: proofs
% Section: inference-patterns

\documentclass[../../../include/open-logic-section]{subfiles}

\begin{document}

\olfileid{mth}{prf}{pat}

\olsection{أنماط الاستدلال}

تتألف البراهين من استدلالات مفردة. وعندما نجري استدلالًا، نشير إلى
ذلك عادة باستعمال كلمة مثل «إذن»، أو «وعليه»، أو «لذلك». وكثيرًا ما
يعتمد الاستدلال على حقيقة أو حقيقتين متاحتين لنا سلفًا في برهاننا---
فقد تكون شيئًا افترضناه، أو شيئًا استنتجناه بالفعل باستدلال. ولزيادة
الوضوح، قد نضع علامات على هذه الأشياء، ونشير في الاستدلال إلى
العبارات الأخرى التي نستخدمها فيه. وكثيرًا ما يتضمن الاستدلال أيضًا
شرحًا لـ\emph{سبب} لزوم النتيجة الجديدة من الأشياء التي سبقتها. وثمة
أنماط استدلال شائعة تُستخدم كثيرًا جدًا في البراهين؛ وسنستعرض بعضها
أدناه. وبعض أنماط الاستدلال، مثل البراهين بالاستقراء، أعقد (وسنناقشها
لاحقًا).

لقد ناقشنا بالفعل نمطًا واحدًا من الاستدلال: فك تعريف، أو تطبيقه.
وعندما نفك تعريفًا، فإننا نعيد ببساطة صياغة شيء يتضمن المعرَّف
باستخدام المُعرِّف. افترض مثلًا أننا أثبتنا بالفعل، في أثناء برهان،
أن~$D = E$ (أ). عندئذ يجوز لنا تطبيق تعريف~$=$ للمجموعات والاستدلال:
«إذن، بالتعريف من (أ)، كل !!{element} من~$D$ هو !!a{element} من~$E$
والعكس صحيح».

ومما يوقع في شيء من الالتباس أننا لا نكتب تسويغ الاستدلال، في كثير
من الأحيان، عند إجرائه فعلًا، بل قبله. افترض أننا لم نثبت بعد
$D = E$، لكننا نريد إثباتها. وإذا كانت~$D = E$ هي النتيجة التي نسعى
إليها، فيمكننا إعادة صياغة هذه الغاية بتطبيق التعريف أيضًا: لإثبات
$D = E$، علينا أن نثبت أن كل !!{element} من~$D$ هو !!a{element} من
$E$، والعكس صحيح. ومن ثم يكون برهاننا على الصورة الآتية: (أ) أثبت أن
كل !!{element} من~$D$ هو !!a{element} من~$E$؛ (ب) كل !!{element} من
$E$ هو !!a{element} من~$D$؛ (ج) إذن، من (أ) و(ب) وبتعريف~$=$،
$D = E$. لكننا لا نكتب البرهان عادة بهذه الطريقة. بل قد نكتب شيئًا
مثل:
\begin{quote}
نريد أن نبيّن~$D = E$. وبحسب تعريف~$=$، يساوي هذا بيان أن كل
!!{element} من~$D$ هو !!a{element} من~$E$، والعكس صحيح.

(أ) \dots{} (برهان أن كل !!{element} من~$D$ هو !!a{element} من~$E$)
\dots

(ب) \dots{} (برهان أن كل !!{element} من~$E$ هو !!a{element} من~$D$)
\dots
\end{quote}

\subsection{استخدام اقتران}

لعل أبسط نمط استدلال هو استخلاص أحد مقترني اقتران بوصفه نتيجة.
وبعبارة أخرى: إذا افترضنا أو أثبتنا بالفعل~$p$ و$q$، حق لنا أن نستنتج
$p$ (وأن نستنتج~$q$ أيضًا). وهذا استدلال أساسي إلى حد أنه لا يُذكر
كثيرًا. فمتى فككنا تعريف~$D = E$ مثلًا، نكون قد أثبتنا أن كل
!!{element} من~$D$ هو !!a{element} من~$E$، والعكس صحيح. ومن ذلك
يمكننا أن نستنتج أن كل !!{element} من~$E$ هو !!a{element} من~$D$
(وهذا هو جزء «والعكس صحيح»).

\subsection{إثبات اقتران}

يكون ما يُطلب منك إثباته أحيانًا على صورة اقتران؛ فيُطلب منك أن
«تثبت~$p$ و$q$». وفي هذه الحالة، ليس عليك إلا فعل شيئين: أن تثبت
$p$، ثم تثبت~$q$. ويمكنك تقسيم برهانك إلى قسمين وتسميتهما لزيادة
الوضوح. وعند كتابة ملاحظاتك الأولى، قد تكتب «(1) أثبت~$p$» في أعلى
الصفحة، و«(2) أثبت~$q$» في منتصفها. (وقد لا يُطلب منك صراحة، بالطبع،
إثبات اقتران، لكنك تجد أن برهانك يتطلب إثبات اقتران. فإذا طُلب منك
مثلًا إثبات~$D = E$، ستجد بعد فك تعريف~$=$ أن عليك إثبات: كل
!!{element} من~$D$ هو !!a{element} من~$E$ \emph{و}كل !!{element} من
$E$ هو !!a{element} من~$D$.)

\subsection{إثبات فصل}

عندما يكون ما تثبته على صورة فصل (أي إنه عبارة من الصورة «$p$ أو
$q$»)، يكفي أن تبيّن صدق أحد المفصولين. لكن لا يحدث في الأساس قط أن
يلزم أحد المفصولين مباشرة من افتراضات مبرهنتك. والأكثر شيوعًا أن تكون
افتراضات مبرهنتك نفسها فصلية، أو أن تكون بصدد بيان أن جميع الأشياء
من ضرب معين لها إحدى خاصيتين، لكن بعض الأشياء له الخاصية الأولى
وبعضها الآخر له الثانية. وهنا يفيد البرهان بالحالات (انظر أدناه).

\subsection{البرهان الشرطي}

تأتي مبرهنات كثيرة ستصادفها على صورة شرطية (أي طُلب منك بيان أنه إذا
تحققت~$p$، صدقت~$q$ أيضًا). وإعداد هذه الحالات لطيف وسهل---افترض
ببساطة مقدم الشرطية (أي~$p$ في هذه الحالة)، وأثبت منه التالي~$q$.
فإذا كانت مبرهنتك تقول «إذا~$p$ فإن~$q$»، تبدأ برهانك بعبارة
«افترض~$p$»، وينبغي أن تكون قد أثبتت~$q$ في نهايته.

وقد تُصاغ الشرطيات بطرائق مختلفة. فبدلًا من «إذا~$p$ فإن~$q$»، قد
تقول مبرهنة إن «$p$ فقط إذا~$q$»، أو «$q$ إذا~$p$»، أو «$q$، بشرط
$p$». وكلها تعني الشيء نفسه وتتطلب افتراض~$p$ وإثبات~$q$ من ذلك
الافتراض. وتذكر أن التكافؤ («$p$ إذا وفقط إذا~$q$») ليس في الحقيقة
إلا شرطيتين مجموعتين: إذا~$p$ فإن~$q$، وإذا~$q$ فإن~$p$. ومن ثم لا
يلزمك إلا إجراء البرهان الشرطي مرتين: مرة للشرطية الأولى وأخرى
للثانية. لكن يمكن أحيانًا إثبات عبارة «إذا وفقط إذا» بوصل مجموعة من
عبارات «إذا وفقط إذا» الأخرى بحيث تبدأ بـ«$p$» وتنتهي بـ«$q$»---لكن
عليك عندئذ التأكد من أن كل خطوة هي حقًا «إذا وفقط إذا».

\subsection{الادعاءات الكلية}

استخدام ادعاء كلي بسيط: إذا صدق شيء على أي شيء، صدق على كل شيء معين.
فإذا كانت فرضية برهانك مثلًا~$A \subseteq B$، فهذا يعني (بفك تعريف
$\subseteq$) أنه لكل~$x \in A$ يكون~$x \in B$. ولذلك، إذا كنت تعلم
بالفعل أن~$z \in A$، أمكنك استنتاج~$z \in B$.

وقد يبدو إثبات ادعاء كلي عسيرًا قليلًا. وتأخذ هذه العبارات عادة
الصورة الآتية: «إذا كانت لـ$x$ الخاصية~$P$، فله الخاصية~$Q$»، أو «كل
$P$ هو~$Q$». وقد لا تطابق هذه الصورة تمامًا بالطبع، وتحتاج إلى شيء
من المران لتعرف المطلوب إثباته على وجه الدقة. لكننا كثيرًا ما نضطر
إلى إثبات أن جميع الأشياء ذات خاصية ما لها خاصية معينة أخرى.

وطريقة إثبات ادعاء كلي هي إدخال أسماء أو متغيرات للأشياء التي لها
الخاصية الأولى، ثم بيان أن لها الخاصية الأخرى أيضًا. وقد نصوغ ذلك
بالقول إنه لكي تثبت شيئًا لـ\emph{كل}~$P$، عليك إثباته لـ$P$
\emph{اعتباطي}. والاسم المدخل اسم لـ$P$ اعتباطي. ونستخدم عادة حروفًا
مفردة أسماء لهذه الأشياء الاعتباطية، وتتبع الحروف في العادة أعرافًا:
فنستخدم مثلًا~$n$ للأعداد الطبيعية، و$!A$ لـ!!{formula}s، و$A$
للمجموعات، و$f$ للدوال، وهلم جرًا.

والحيلة هي الحفاظ على العمومية طوال البرهان. تبدأ بافتراض أن شيئًا
اعتباطيًا («$x$») له الخاصية~$P$، وتبيّن (استنادًا فقط إلى التعريفات
أو إلى ما يجوز لك افتراضه) أن لـ$x$ الخاصية~$Q$. ولأنك لم تحدد ما
$x$ على وجه الخصوص سوى أن له الخاصية~$P$، يمكنك أن تقرر أن كل شيء
له~$P$ يملك الخاصية~$Q$. وباختصار، ينوب~$x$ عن \emph{جميع} الأشياء
ذات الخاصية~$P$.

\begin{prop}
  لكل مجموعتين~$A$ و$B$، يكون~$A \subseteq A \cup B$.
\end{prop}

\begin{proof}
  لتكن~$A$ و$B$ مجموعتين اعتباطيتين. نريد أن نبيّن
  $A \subseteq A \cup B$. وبحسب تعريف~$\subseteq$، يساوي هذا: لكل
  $x$، إذا كان~$x \in A$، فإن~$x \in A \cup B$. فليكن
  $x \in A$ !!{element} اعتباطيًا من~$A$. علينا أن نبيّن
  $x \in A \cup B$. ولأن~$x \in A$، يكون~$x \in A$ أو~$x \in B$.
  ومن ثم~$x \in \Setabs{x}{x \in A \lor x \in B}$. لكن هذا يعني،
  بحسب تعريف~$\cup$، أن~$x \in A \cup B$.
\end{proof}


\subsection{البرهان بالحالات}

افترض أن لديك فصلًا بوصفه افتراضًا أو نتيجة مثبتة بالفعل---أي إنك
افترضت أو أثبتت صدق~$p$ أو~$q$. وتريد إثبات~$r$. وتفعل ذلك في
خطوتين: تفترض أولًا صدق~$p$ وتثبت~$r$، ثم تفترض صدق~$q$ وتثبت~$r$
مرة أخرى. وينجح ذلك لأننا نفترض أو نعلم أن أحد البديلين متحقق.
وتثبت الخطوتان أن أيًا منهما يكفي لصدق~$r$. (وإذا صدقا معًا، لم يكن
لدينا سبب واحد، بل سببان، لصدق~$r$. ولا يلزم أن نثبت على حدة صدق~$r$
بافتراض~$p$ و$q$ معًا.) ولكي نشير إلى ما نفعله، نعلن أننا «نميز
الحالات». افترض مثلًا أننا نعلم~$x \in B \cup C$. وتُعرّف
$B \cup C$ بأنها~$\Setabs{x}{x \in B \text{ أو } x \in C}$. أي
إن~$x \in B$ أو~$x \in C$ بحسب التعريف. ونثبت من ذلك~$x \in A$
بافتراض~$x \in B$ أولًا وإثبات~$x \in A$ من هذا الافتراض، ثم
بافتراض~$x \in C$ وإثبات~$x \in A$ مرة أخرى من هذا الافتراض. وتكتب
«نميز الحالات» تحت الافتراض، ثم «الحالة (1):~$x \in B$» تحتها،
و«الحالة (2):~$x \in C$» في منتصف الصفحة. ثم تملأ النصف الأعلى
والنصف الأسفل من الصفحة.

يفيد البرهان بالحالات بوجه خاص إذا كان ما تثبته نفسه فصليًا. وفيما
يأتي مثال بسيط:

\begin{prop}
افترض~$B \subseteq D$ و$C \subseteq E$. عندئذ
$B \cup C \subseteq D \cup E$.
\end{prop}

\begin{proof}
  افترض (أ)~$B \subseteq D$ و(ب)~$C \subseteq E$. بحسب التعريف، كل
  $x \in B$ يحقق أيضًا~$x \in D$ (ج)، وكل~$x \in C$ يحقق أيضًا
  $x \in E$ (د). ولكي نبيّن~$B \cup C \subseteq D \cup E$، علينا أن
  نبيّن أنه إذا كان~$x \in B \cup C$، فإن~$x \in D \cup E$ (بحسب
  تعريف~$\subseteq$). ويكون~$x \in B \cup C$ إذا وفقط إذا كان
  $x \in B$ أو~$x \in C$ (بحسب تعريف~$\cup$). وبالمثل، يكون
  $x \in D \cup E$ إذا وفقط إذا كان~$x \in D$ أو~$x \in E$. ومن ثم
  علينا أن نبيّن: لكل~$x$، إذا كان~$x \in B$ أو~$x \in C$، فإن
  $x \in D$ أو~$x \in E$.

  \begin{quote}
  لم نفعل حتى الآن إلا فك التعريفات!{} فقد أعدنا صياغة قضيتنا من دون
  $\subseteq$ و$\cup$، وبقي علينا أن نحاول إثبات ادعاء شرطي كلي.
  وبحسب ما ناقشناه أعلاه، يجري ذلك بافتراض أن~$x$ شيء نفترض صدق جزء
  «إذا» عليه، ثم نواصل فنبين صدق جزء «فإن» أيضًا. وبعبارة أخرى،
  سنفترض~$x \in B$ أو~$x \in C$ ونبيّن~$x \in D$ أو
  $x \in E$.\footnote{لا تفعل هذه الفقرة إلا شرح ما نفعله---فهي ليست
  جزءًا من البرهان، ولا يلزمك الخوض في كل هذه التفاصيل عندما تكتب
  براهينك.}
  \end{quote}

  افترض~$x \in B$ أو~$x \in C$. وعلينا أن نبيّن~$x \in D$ أو
  $x \in E$. نميز الحالات.

  الحالة 1:~$x \in B$. بموجب (ج)،~$x \in D$. ومن ثم~$x \in D$ أو
  $x \in E$. (لقد أجرينا هنا الاستدلال المناقش في القسم الفرعي
  السابق!)

  الحالة 2:~$x \in C$. بموجب (د)،~$x \in E$. ومن ثم~$x \in D$ أو
  $x \in E$.
 \end{proof}


\subsection{إثبات ادعاء وجودي}

عندما يُطلب منك إثبات ادعاء وجودي، يكون السؤال في العادة على الصورة
«أثبت وجود~$x$ بحيث~$\dots x \dots$»، أي إن شيئًا ما له الخاصية
الموصوفة بـ«$\dots x \dots$». وفي هذه الحالة، عليك تعيين شيء مناسب
وبيان أن له الخاصية المطلوبة. ويبدو ذلك مباشرًا، لكن برهانًا من هذا
النوع قد يكون عسيرًا. فهو يتضمن عادة \emph{بناء} شيء أو \emph{تعريفه}
وإثبات أن للشيء المعرَّف هكذا الخاصية المطلوبة. وقد يصعب العثور على
الشيء الصحيح، وقد يصعب إثبات أن له الخاصية المطلوبة، وأحيانًا يصعب
حتى بيان أنك نجحت أصلًا في تعريف شيء!{}

وعادة ما تكتب ذلك بتحديد الشيء، مثل «ليكن~$x$ هو \dots» (حيث تحدد
\dots{} الشيء الذي تقصده)، وقد تثبت أن~$\dots$ تصف فعلًا شيئًا
موجودًا، ثم تواصل فتبين أن لـ$x$ الخاصية~$Q$. وفيما يأتي مثال بسيط.

\begin{prop}
  افترض~$x \in B$. عندئذ توجد~$A$ بحيث~$A \subseteq B$ و
  $A \neq \emptyset$.
\end{prop}

\begin{proof}
  افترض~$x \in B$. لتكن~$A = \{x\}$.
  \begin{quote}
    عرّفنا هنا المجموعة~$A$ بتعداد !!{element}sها. ولأننا نفترض أن
    $x$ شيء، ولأن بوسعنا دائمًا تكوين مجموعة بتعداد !!{element}sها،
    فلا يلزم أن نبيّن هنا أننا نجحنا في تعريف مجموعة~$A$. لكن ما زال
    علينا أن نبيّن أن لـ$A$ الخواص التي تتطلبها القضية. ولا يكتمل
    البرهان من دون ذلك!{}
  \end{quote}
  بما أن~$x \in A$، يكون~$A \neq \emptyset$.
  \begin{quote}
    يعتمد هذا على تعريف~$A$ بأنها~$\{x\}$، وعلى الحقيقتين الواضحتين
    $x \in \{x\}$ و$x \notin \emptyset$.
  \end{quote}
  لأن~$x$ هو !!{element} الوحيد من~$\{x\}$، ولأن~$x \in B$، يكون كل
  !!{element} من~$A$ أيضًا !!a{element} من~$B$. وبحسب تعريف
  $\subseteq$، يكون~$A \subseteq B$.
\end{proof}

\subsection{استخدام الادعاءات الوجودية}

افترض أنك تعرف صدق ادعاء وجودي ما (لأنك أثبته، أو لأنه فرضية يجوز لك
استخدامها)، وليكن «لبعض~$x$،~$x \in A$»، أو «يوجد~$x \in A$». فإذا
أردت استخدامه في برهانك، أمكنك ببساطة أن تتظاهر بأن لديك اسمًا لأحد
الأشياء التي تقول فرضيتك إنها موجودة. ولأن~$A$ تحتوي شيئًا واحدًا على
الأقل، توجد أشياء قد يشير إليها ذلك الاسم. وقد لا تستطيع بالطبع
تعيين أحدها أو وصفه أكثر (سوى أنه~$\in A$). لكنك تستطيع، لغرض
البرهان، أن تتظاهر بأنك عينته وأن تسميه. ومن المهم أن تختار اسمًا لم
تستخدمه بالفعل (ولم يرد في فرضياتك)، وإلا فقد تختل الأمور. وتشير إلى
ذلك في برهانك بالانتقال من «لبعض~$x$،~$x \in A$» إلى «ليكن
$a \in A$». ويمكنك الآن أن تستدل بشأن~$a$، وتستخدم فرضيات أخرى، وهلم
جرًا، إلى أن تبلغ نتيجة~$p$. فإذا لم تعد~$p$ تذكر~$a$، كانت~$p$
مستقلة عن الافتراض~$a \in A$، وتكون قد بينت أنها تلزم فقط من الافتراض
«لبعض~$x$،~$x \in A$».

\begin{prop}
إذا كان~$A \neq \emptyset$، فإن~$A \cup B \neq \emptyset$.
\end{prop}

\begin{proof}
  افترض~$A \neq \emptyset$. ومن ثم، لبعض~$x$، يكون~$x \in A$.
  \begin{quote}
    أعدنا هنا أولًا صياغة فرضية القضية فحسب. تخفي هذه الفرضية، أي
    $A \neq \emptyset$، ادعاء وجوديًا لا تصل إليه إلا بفك عدة
    تعريفات. يخبرنا تعريف~$=$ أن~$A = \emptyset$ إذا وفقط إذا كان،
    لكل~$x$، إذا~$x \in A$ فإن~$x \in \emptyset$، وإذا
    $x \in \emptyset$ فإن~$x \in A$. وبنفي الطرفين، نحصل على أن
    $A \neq \emptyset$ إذا وفقط إذا وُجد~$x \in A$ بحيث
    $x \notin \emptyset$، أو وُجد~$x \in \emptyset$ بحيث
    $x \notin A$. ولأن العبارة~$x \in \emptyset$ لا تصدق على أي~$x$،
    لا يمكن أن يصدق المفصول الثاني قط، كما يختزل «$x \in A$ و
    $x \notin \emptyset$» إلى~$x \in A$ فحسب. وهكذا
    $A \neq \emptyset$ إذا وفقط إذا كان~$x \in A$ لبعض~$x$. وهذا
    ادعاء وجودي. ونستخدم الآن ذلك الادعاء الوجودي بإدخال اسم لأحد
    !!{element}s~$A$:
  \end{quote}
  ليكن~$a \in A$.
  \begin{quote}
    أدخلنا الآن اسمًا لأحد الأشياء التي~$\in A$. وسنواصل الاستدلال
    بشأن~$a$، لكننا سنحرص على ألا نفترض إلا~$a \in A$:
  \end{quote}
  لأن~$a \in A$، يكون~$a \in A \cup B$ بحسب تعريف~$\cup$. ومن ثم،
  لبعض~$x$، يكون~$x \in A \cup B$، أي~$A \cup B \neq \emptyset$.
  \begin{quote}
    انتقلنا في الخطوة الأخيرة من «$a \in A \cup B$» إلى «لبعض~$x$،
    $x \in A \cup B$». ولم تعد هذه العبارة تذكر~$a$، ولذلك نعلم أن
    «لبعض~$x$،~$x \in A \cup B$» تلزم من «لبعض~$x$،~$x \in A$»
    وحدها. لكن ذلك يعني~$A \cup B \neq \emptyset$.
  \end{quote}
\end{proof}

قد يكون من حسن الممارسة أن تبقي المتغيرات المقيدة مثل «$x$» منفصلة
عن الأسماء الافتراضية مثل~$a$، كما فعلنا. لكننا في الممارسة كثيرًا
ما لا نفعل ذلك، ونستخدم~$x$ فحسب، هكذا:
\begin{quote}
افترض~$A \neq \emptyset$، أي يوجد~$x \in A$. وبحسب تعريف~$\cup$،
$x \in A \cup B$. ومن ثم~$A \cup B \neq \emptyset$.
\end{quote}
لكن عندما تفعل ذلك، يجب أن تحرص حرصًا زائدًا على استخدام رموز~$x$
و$y$ مختلفة لادعاءات وجودية مختلفة. فالآتي مثلًا \emph{ليس} برهانًا
صحيحًا على «إذا كان~$A \neq \emptyset$ و$B \neq \emptyset$، فإن
$A \cap B \neq \emptyset$» (وهي عبارة غير صادقة).
\begin{quote}
افترض~$A \neq \emptyset$ و$B \neq \emptyset$. ومن ثم لبعض~$x$ يكون
$x \in A$، ولبعض~$x$ أيضًا يكون~$x \in B$. ولأن~$x \in A$ و
$x \in B$، يكون~$x \in A \cap B$ بحسب تعريف~$\cap$. ومن ثم
$A \cap B \neq \emptyset$.
\end{quote}
هل تستطيع تعيين موضع الخطوة غير الصحيحة وشرح سبب عدم ثبوت النتيجة؟

\end{document}
